\documentclass{article}
\usepackage[utf8]{inputenc}
\usepackage[T1]{fontenc}
\usepackage[a4paper,total={6in,8in}]{geometry}
\usepackage{amsmath,amssymb,amsfonts}
\usepackage{graphicx}
\usepackage{hyperref}
\usepackage{enumitem}
\usepackage{listings}
\usepackage{xcolor}
\DeclareMathOperator{\Var}{Var}
\DeclareMathOperator{\Cov}{Cov}
\providecommand{\tightlist}{\setlength{\itemsep}{0pt}\setlength{\parskip}{0pt}}
\lstset{basicstyle=\ttfamily\scriptsize,breaklines=true,columns=fullflexible,
  keepspaces=true,showstringspaces=false,frame=single}
\hypersetup{colorlinks=true,urlcolor=blue,linkcolor=blue}
\setlength{\emergencystretch}{2em}

\title{Math 126 Problem 14: Visual Applications of SVD}
\author{} % Put your name here.
\date{Spring 2027}
\begin{document}
\maketitle

\hypertarget{submission-instructions}{%
\subsection*{Submission Instructions}\label{submission-instructions}}

Please upload your work to Gradescope by \textbf{11:59pm on Monday, March 29, 2027 (Pacific time).}

\begin{itemize}
\item
  You are encouraged to discuss ideas and work with your classmates. However, you \emph{must acknowledge} your collaborators at the top of each solution on which you collaborated with others and you \emph{must write} your solutions independently.
\item
  Your solutions to theory questions must be written legibly, or typeset in LaTeX or markdown.
\item
  I recommend that you write your solutions in a Python notebook using Google Colab.
\item
  You should submit your solutions as a \emph{single PDF} via the assignment on Gradescope.
\end{itemize}

\textbf{Grading:} The point of the problem set is for \emph{you} to learn. To this end, I hope to disincentivize the use of LLMs by \emph{not} grading your work for correctness. Instead, you will grade your own work by comparing it to my solutions. This self-grade is due the Friday \emph{after} the problem set is due, also on Gradescope.

\noindent Original notebook: \href{https://colab.research.google.com/drive/131XUvWaOZQLDnd_PZGU1o8yenD8XoOgP?usp=sharing}{Google Colab}.

% Source notebook cell 2
\hypertarget{visual-applications-of-svd}{%
\subsection*{Visual Applications of SVD}\label{visual-applications-of-svd}}

\hypertarget{part-a-low-rank-approximations-to-images}{%
\subsection*{Part A: Low Rank Approximations to Images}\label{part-a-low-rank-approximations-to-images}}

Load an image of your choice from the internet as a matrix \(\mathbf{X} \in \mathbb{R}^{n \times d}\). Use \texttt{numpy}'s built-in tools to compute its SVD. Plot the singular values by their index; this is known as a \emph{spectrum} plot.

Next, find plot the best rank-\(k\) approximation for \(k=1, 5, 10, 50, 100, 500, 1000\). How do the rank-\(k\) approximations relate to the spectrum?

\hypertarget{part-b-anomaly-detection}{%
\subsection*{Part B: Anomaly Detection}\label{part-b-anomaly-detection}}

In this problem, we'll consider a different way of building the matrix \(\mathbf{X}\). Download a surveillance video (e.g., from \href{https://raw.githubusercontent.com/AarohiSingla/Speed-detection-of-vehicles/main/highway.mp4}{here}). Each row of \(\mathbf{X}\) is a time step where the vector is matrix representation of the corresponding image \emph{flattened}. Note: Please compress the image first so the flattened image is a reasonable size.

Next, compute the SVD of your matrix \(\mathbf{X}\). Try subtracting the right singular vector from a row of \(\mathbf{X}\) with abnormal movement. What do you see?

\end{document}
